\documentclass[11pt]{article}
\usepackage[T1]{fontenc}
\usepackage{lmodern}
\usepackage[margin=1in]{geometry}
\usepackage{amsmath,amssymb,amsthm,mathtools}
\usepackage{enumitem}
\usepackage{microtype}
\usepackage{xurl}
\usepackage[colorlinks=true,linkcolor=blue,citecolor=blue,urlcolor=blue]{hyperref}
\newtheorem{theorem}{Theorem}[section]
\newtheorem{proposition}[theorem]{Proposition}
\newtheorem{lemma}[theorem]{Lemma}
\theoremstyle{remark}
\newtheorem{remark}[theorem]{Remark}
\newtheorem{example}[theorem]{Example}
\newcommand{\D}{\mathbb D}
\newcommand{\Sone}{\mathbb S^1}
\newcommand{\Cpair}{\mathcal C}
\newcommand{\BA}{\mathcal B_A}
\newcommand{\Fset}{\mathcal F}
\newcommand{\Kset}{\mathcal K}
\newcommand{\Varea}{\mathcal V}
\newcommand{\dd}{\,d}
\DeclareMathOperator{\Area}{Area}
\DeclareMathOperator{\supp}{supp}
\hypersetup{pdftitle={A variational characterization of geodesic scattering on simple disks},pdfauthor={Kaiwen Jin}}
\title{A variational characterization of geodesic scattering on simple disks}
\author{Kaiwen Jin\\\small\href{mailto:yui.ui.siki@gmail.com}{yui.ui.siki@gmail.com}}
\date{October 12, 2026}
\begin{document}
\maketitle
\begin{abstract}
We give necessary and sufficient conditions for a boundary map to be the
geodesic scattering map of a smooth simple Riemannian metric on a fixed disk
with a prescribed boundary metric. The conditions depend only on the proposed
map and that boundary metric. Boundary invariant functions first determine a
normalized conformal boundary coordinate. In the resulting Euclidean disk,
length inequalities define a nonempty convex set of densities whose strong
$L^2$ closure has a unique element of least norm. Realizability is equivalent
to smooth positivity and the prescribed boundary value of this element,
explicit geodesic and Jacobi equation tests, and a scalar area equality.
The proof transfers the length inequalities through the closure by integrating
over the geodesic family. Strict monotonicity of the area functional then
forces equality of boundary distances and hence of scattering maps. The
characterization is exact and qualitative; no finite decision procedure or
stability estimate is asserted.
\end{abstract}

\section{Introduction}\label{sec:intro}
The scattering map records the point and direction at which a geodesic
leaves a manifold after entering at its boundary. On a simple surface,
every inward direction has a finite flight and distinct boundary endpoints.
For a fixed boundary metric, the tangential components of the incoming and
outgoing unit covectors give a common space on which to compare these maps.
The range question asks which proposed maps are produced by smooth simple
Riemannian metrics. It is a consistency question for the directions of rays,
before one attempts to recover an interior metric.

This question differs from uniqueness for data already known to come from
a metric. Pestov and Uhlmann prove that two-dimensional compact simple
Riemannian manifolds are boundary distance rigid, and their
Theorem~1.3 recovers the Dirichlet-to-Neumann map from actual scattering
data \cite{PU}. The metric-to-scattering range question is posed separately
in \cite[Section~15.1, problem~2]{PSU}. Bohr and Paternain's
Theorem~1.5 characterizes non-Abelian scattering data on an already fixed
simple surface \cite{BP}; the geometric scattering relation of that surface
is an input there. Lens-map realization results for Finsler metrics
\cite{BI} also concern a different class of interior geometries.
The problem considered here allows an arbitrary proposed boundary map and
requires a Riemannian realization. The geometric and transport results
used below are those developed in \cite{PSU}.

Fix a smooth compact disk $M$ and a smooth positive boundary metric $h$.
The requested characterization is in terms of $A$ and $h$ alone, with
no additional interior metric among the data. We give such a characterization
by selecting a density before testing realizability.

\medskip
\noindent\textbf{Main result, informal form.}
\emph{A map on the entire open unit cotangent bundle of the boundary is
realizable if and only if it passes the boundary action and boundary
Fourier tests below, and the uniquely selected least-norm density is smooth
and positive with the prescribed boundary value, passes the fixed
geodesic and Jacobi equation tests, and has the area determined by the
boundary action. There is no smallness, genericity, or local perturbation
assumption.}
\medskip

The boundary Fourier test extracts the traces of holomorphic functions
for actual scattering data, and a three-point normalization selects a
circle coordinate. One then considers all positive smooth densities on
a fixed Euclidean disk whose lengths dominate the proposed boundary
distances. This convex set is always nonempty after the boundary tests,
even for data that have no realization. Its strong $L^2$ closure therefore
selects one density class without making a metric-existence assumption.
The remaining conditions test this fixed class rather than search among
simple metrics.

The main analytical difficulty is that restriction to an individual curve
is not continuous in $L^2$ on the disk. Santal\'o's formula instead controls
the error integrated over all geodesics of the selected smooth metric.
This yields a pointwise lower bound on its boundary distance, by continuity.
A symmetric integration-by-parts identity shows that the action area is
strictly increasing under a strict increase of boundary distance with
positive mixed derivative. Thus one scalar area equality forces exact
equality of all boundary distances. First variation then identifies the
scattering map.

The result is a qualitative characterization of the full range. Its tests
include infinitely many curves, smoothness of an $L^2$ minimizer, and exact
ODE and scalar conditions. It does not supply a finite computation or
quantitative stability, and it does not address non-simple surfaces or
Finsler realization.

Section~\ref{sec:conditions} states the conditions and theorem.
Section~\ref{sec:boundary} proves the boundary and conformal facts for
actual simple metrics. Section~\ref{sec:variation} proves the universal
convex construction and the two directions of the characterization.
Throughout, the exit covector is recorded without reversal, Fourier
averages have the factor $1/(2\pi)$, and all $L^2$ spaces and closures of
densities are real and use Euclidean area.

\section{Boundary tests and the selected density}\label{sec:conditions}
Choose an orientation, an $h$-arclength origin, and three distinct
positively cyclically ordered boundary points $s_1,s_2,s_3$.
Write $L>0$ for the $h$-length of the boundary and identify it with
$\mathbb R/L\mathbb Z$. The positive unit tangent $T$ is chosen so that
$(T,\nu)$ is positive when $\nu$ is the inward unit normal.
Then
\[
 Z=(\mathbb R/L\mathbb Z)\times(-1,1)
\]
identifies with $B_h^*\partial M$: $p$ denotes the coefficient of the
tangential covector $p\,ds$. Start with an arbitrary set map $A:Z\to Z$.
The following conditions are evaluated in order. If a condition fails,
the map is rejected and later constructions are not needed.

\subsection{The boundary action: condition (C1)}\label{subsec:c1}
Require that $A$ be a smooth diffeomorphism. Write $A(s,p)=(t,q)$ and
require the endpoint map
\[
 E:Z\longrightarrow\Cpair,
 \qquad E(s,p)=(s,t),\qquad
 \Cpair=\{(s,t)\in(\mathbb R/L\mathbb Z)^2:s\ne t\},
\]
to be a smooth diffeomorphism. Its inverse defines functions
$P(s,t)=p$ and $Q(s,t)=q$. For the unique forward lift $s<t<s+L$, put
\begin{equation}\label{eq:primitive}
 D(s,t)=\int_s^t Q(s,u)\dd u.
\end{equation}
The integral converges at its lower endpoint since $|Q|<1$.
Its profile $d(s,\delta)=D(s,s+\delta)$, initially $0<\delta<L$,
must extend smoothly to $(\mathbb R/L\mathbb Z)\times[0,L]$ and satisfy
\begin{gather}\label{eq:profile}
 d(s,0)=d(s,L)=0,\qquad
 d_\delta(s,0)=1,\qquad d_\delta(s,L)=-1,\\
 D>0,\qquad D(s,t)=D(t,s),\qquad
 D_s=-P,\qquad D_t=Q,\qquad D_{st}>0\quad\hbox{on }\Cpair.
 \label{eq:action}
\end{gather}
Derivatives of $D$ treat its endpoints as independent variables; hence
$D_s=d_s-d_\delta$ and $D_{st}=d_{s\delta}-d_{\delta\delta}$.
Define the action area
\begin{equation}\label{eq:area-action}
 \Varea(D)=\frac1{2\pi}\int_{\Cpair}D D_{st}\dd s\dd t.
\end{equation}
This is a finite integral of the smooth profile on the compact rectangle,
using positive product arclength measure $ds\,dt=ds\,d\delta$.

\subsection{A conformal boundary coordinate: condition (C2)}\label{subsec:c2}
For a metric inducing $h$, a boundary unit vector has the form
$\cos\theta\,T+\sin\theta\,\nu$.
The full boundary torus is the abstract space
$S=(\mathbb R/L\mathbb Z)\times(\mathbb R/2\pi\mathbb Z)$.
Define a piecewise map on it by
\begin{equation}\label{eq:full-lift}
 \alpha_A(s,\theta)=
 \begin{cases}
 (t,2\pi-\arccos q), &0<\theta<\pi,
       \quad A(s,\cos\theta)=(t,q),\\
 (u,\arccos p), &\pi<\theta<2\pi,
       \quad A^{-1}(s,\cos\theta)=(u,p),\\
 (s,\theta), &\theta=0,\pi.
 \end{cases}
\end{equation}
Require that $\alpha_A$ be smooth on $S$. Its square is the identity by
construction, so it is then a smooth diffeomorphism.
For $U\in C^\infty(S;\mathbb C)$ use the coefficients
\[
 U_k(s)=\frac1{2\pi}\int_0^{2\pi}U(s,\theta)e^{-ik\theta}\dd\theta.
\]
Let $\BA$ consist of the functions $U_0$ obtained from such $U$ with
\begin{equation}\label{eq:boundary-fourier}
 U\circ\alpha_A=U,\qquad
 U(s,\theta+\pi)=U(s,\theta),\qquad
 U_k=0\quad(k<0).
\end{equation}
This definition uses only functions on the boundary torus.

Set $\omega=e^{2\pi i/3}$. Require that there be exactly one smooth
orientation-preserving circle diffeomorphism
\begin{equation}\label{eq:phi}
 \phi:\mathbb R/L\mathbb Z\longrightarrow\Sone,
 \qquad \phi\in\BA,\qquad \phi(s_j)=\omega^{j-1}
       \quad(j=1,2,3).
\end{equation}
Here membership in $\BA$ treats $\phi$ as a complex-valued function.
On the unit circle define the positive boundary density $b$ and action
$D_\phi$ by
\begin{equation}\label{eq:boundary-density}
 b(\phi(s))=\frac1{|\phi_s(s)|},\qquad
 D_\phi(a,c)=D(\phi^{-1}(a),\phi^{-1}(c))\quad(a\ne c).
\end{equation}
The angular derivative of an orientation-preserving diffeomorphism is
positive, so the denominator never vanishes.

\subsection{The convex construction: condition (C3)}\label{subsec:c3}
All further constructions take place in the fixed closed Euclidean unit
disk $\overline\D$. Smooth functions on this disk mean restrictions of
smooth functions on a neighborhood. Let
\[
 H=L^2(\D,dx\,dy;\mathbb R),\qquad
 \|u\|_H^2=\int_\D u^2\dd x\dd y,
\]
with equality of functions understood almost everywhere. Let $\Fset$
be the set of $\sigma\in C^\infty(\overline\D;\mathbb R)$, strictly
positive on $\overline\D$, such that
\begin{equation}\label{eq:all-curves}
 \int_0^1\sigma(\gamma(r))|\gamma'(r)|\dd r
       \ge D_\phi(a,c)
\end{equation}
for every pair $a\ne c$ on $\Sone$ and every piecewise $C^1$ curve
$\gamma:[0,1]\to\overline\D$ joining them. Curves along the boundary,
self-intersections, and repeated segments are included.
No boundary-value, curvature, or simplicity requirement is imposed on
$\sigma$ in this definition. Let $\Kset$ be the strong $H$-closure of
$\Fset$.

Lemma~\ref{lem:projection} shows that $\Fset$ is nonempty and convex for
every map passing (C1)--(C2), and that $\Kset$ has a unique element
$\rho_D$ of smallest $H$ norm. Thus $\rho_D$ is defined before any
realizability test on it. Require that this one $L^2$ class have a smooth
representative $\rho$ that is strictly positive on $\overline\D$, and
require
\begin{equation}\label{eq:rho-trace}
 \rho|_{\Sone}=b.
\end{equation}
A continuous representative, if it exists, is unique, including its
boundary values by continuity. Curve integrals are never evaluated on an
arbitrary representative of an $L^2$ class.

\subsection{Fixed differential equations: condition (C4)}\label{subsec:c4}
If (C3) passes, put $\lambda=\log\rho$, $v(\beta)=(\cos\beta,\sin\beta)$,
$J(x,y)=(-y,x)$ and $K_\rho=-\rho^{-2}\Delta\lambda$.
Require
\begin{equation}\label{eq:convexity}
 1+\partial_r\lambda>0\quad\hbox{on }\Sone,
\end{equation}
where $\partial_r$ is the outward radial derivative. For every interior
initial point and every $\beta\in\mathbb R/2\pi\mathbb Z$, and for every
boundary initial point with $z\cdot v(\beta)<0$, solve the fixed equations
\begin{equation}\label{eq:fixed-ode}
 z'=\frac{v(\beta)}{\rho(z)},\qquad
 \beta'=\frac{\nabla\lambda(z)\cdot Jv(\beta)}{\rho(z)},\qquad
 j''+K_\rho(z)j=0,\quad j(0)=0,\quad j'(0)=1.
\end{equation}
Require a finite first exit time $\tau>0$, with $|z(t)|<1$ for
$0<t<\tau$, and require
\begin{equation}\label{eq:jacobi-positive}
 j(t)>0\quad(0<t\le\tau).
\end{equation}
For an inward boundary state, first exit means first return after
entering the disk. For an interior state it means its first subsequent
encounter with the boundary. The equations are considered only up to
this encounter; a smooth extension of $\rho$ outside the disk does not
change their solution before it. Boundary tangent states have zero flight
time and are excluded.

\subsection{Area saturation and the theorem}\label{subsec:c5}
The final condition (C5) is the exact scalar equality
\begin{equation}\label{eq:saturation}
 \|\rho_D\|_H^2=\Varea(D).
\end{equation}
Neither side is an approximate or numerical quantity.

\begin{theorem}\label{thm:main}
For every smooth compact disk $M$, every smooth positive boundary metric
$h$, and every map $A:B_h^*\partial M\to B_h^*\partial M$, conditions
(C1)--(C5) are necessary and sufficient for $A$ to be the scattering map
of a smooth simple Riemannian metric on $M$ inducing $h$. Their truth
is independent of the initial orientation, arclength origin, and positively
ordered boundary triple.
\end{theorem}

In these conditions, $L$ is the boundary length, $D$ is the primitive
\eqref{eq:primitive}, $\phi$ is the uniquely normalized boundary map
\eqref{eq:phi}, $b$ is the prescribed circle density
\eqref{eq:boundary-density}, and $\rho_D$ is the unique least-norm
element of the strong closure defined by \eqref{eq:all-curves}.
All are determined by $A,h$ and the fixed auxiliary choices. In particular,
the universal construction of $\rho_D$ is distinct from the conclusion
that a realizing metric exists.

\begin{example}\label{ex:euclidean}
For the Euclidean disk of radius $R>0$, with $L=2\pi R$ and marks
$0,L/3,2L/3$, the data are
\[
 A(s,p)=(s+2R\arccos p,p),\qquad
 d(s,\delta)=2R\sin\frac\delta{2R},\qquad
 D_{st}=\frac1{2R}\sin\frac\delta{2R}.
\]
Thus $\Varea(D)=\pi R^2$. The full lift is
$\alpha(s,\theta)=(s+2R\theta,-\theta)$ modulo the two periods.
The function $U=e^{is/R}(1+e^{2i\theta})$ is even, invariant under this
lift, and has modes $0,2$. Its average is $\phi(s)=e^{is/R}$.
Lemma~\ref{lem:coordinate} gives its required uniqueness, and
Proposition~\ref{prop:necessity} identifies $\rho_D=R$.
Here $b=R$, $K_\rho=0$ and $j(t)=t$. Every permitted ray has finite
exit, the expression in \eqref{eq:convexity} is $1$, and all five
conditions hold. This is an exact calculation of the signs and area factor.
\end{example}

\section{Boundary geometry and conformal coordinates}\label{sec:boundary}
We first consider data from an actual smooth simple metric $g$ inducing
$h$. Let $J_g:TM\to TM$ be the positive $g$-orthogonal rotation by
$\pi/2$; in oriented conformal coordinates it agrees with the Euclidean
rotation $J$ in (C4). We use the equivalence between simplicity and strict convexity,
nontrapping, and absence of conjugate points on a disk
\cite[Theorem~3.8.2]{PSU}. Its connecting geodesics minimize length
\cite[Proposition~3.8.4]{PSU}.

\subsection{The action, twist and area}\label{subsec:geometry}
\begin{lemma}\label{lem:geometry}
For actual simple-metric scattering data, (C1) holds and the primitive
$D$ is the boundary distance $D_g$. The full lift \eqref{eq:full-lift}
is smooth. Moreover,
\begin{equation}\label{eq:santalo-area}
 \Area_g(M)=\Varea(D_g).
\end{equation}
\end{lemma}
\begin{proof}
A geodesic joining distinct boundary points enters and exits strictly
transversely. Indeed a boundary defining function positive in the
interior has negative second derivative along a tangent geodesic, by
strict convexity, so that such a geodesic lies outside near the encounter.
The minimizing property also rules out a positive-length return to the
starting point.

Let a connecting geodesic vary smoothly with its endpoints. For its
constant-speed parametrization $c(r,\varepsilon)$, $0\le r\le1$,
write $W=\partial_r c$, $Y=\partial_\varepsilon c$ and length
$\ell(\varepsilon)$. Metric compatibility, torsion-freeness, and integration
by parts give
\[
 \ell'=\ell^{-1}\int_0^1 g(\nabla_rY,W)\dd r
        =g(W/\ell,Y)\big|_{r=1}-g(W/\ell,Y)\big|_{r=0}.
\]
Thus $(D_g)_s=-p$ and $(D_g)_t=q$, with the outgoing covector retained.

To calculate the sign of the mixed derivative, fix the initial boundary
point and vary the inward angle $0<\theta<\pi$. The normal Jacobi field
is $j(t)J_g\dot c(t)$, where $j(0)=0$, $j'(0)=1$ and
$j''+K_g(c(t))j=0$; see \cite[Lemmas~3.7.2, 3.7.6 and 3.7.7]{PSU}.
It is positive until exit, since there are no conjugate points.
Write the exit vector as $qT-r\nu$, $r=\sqrt{1-q^2}>0$.
Differentiating the exit event and pairing with $J_g\dot c(\tau)$ yields
\[
 T\,t_\theta=j(\tau)J_g\dot c(\tau)+\dot c(\tau)\tau_\theta,
 \qquad r t_\theta=j(\tau).
\]
Since $p=\cos\theta$, this gives
\begin{equation}\label{eq:twist}
 t_p=-\frac{j(\tau)}{r\sin\theta},\qquad
 (D_g)_{st}=-p_t=\frac{r\sin\theta}{j(\tau)}>0.
\end{equation}
Unique connecting geodesics give a bijective endpoint map, and this
nonzero derivative gives its smooth inverse. The smooth full scattering
involution, including the tangent states, is
\cite[Lemma~3.3.5]{PSU}; in the $(T,\nu)$ frame it is exactly
\eqref{eq:full-lift}. Its incoming restriction gives a smooth
diffeomorphism on $Z$.

The diagonal behavior is checked in endpoint coordinates, without
requiring smoothness in $p$ at $\pm1$. Extend $g$ to a closed surface
and to a slightly larger simple neighborhood, as in
\cite[Lemma~3.1.8 and Proposition~3.8.7]{PSU}.
The exponential inverse $w(s,\delta)$ of the nearby endpoints
$x(s),x(s+\delta)$ is smooth with $w(s,0)=0$ and
$w_\delta(s,0)=T$. The identity
\[
 \frac{w(s,\delta)}\delta
       =\int_0^1 w_\delta(s,r\delta)\dd r
\]
extends this quotient smoothly and nonvanishingly to $\delta=0$.
For small positive $\delta$ the connecting geodesic lies in $M$ and
\[
 D_g(s,s+\delta)
   =\delta\sqrt{g_{x(s)}(w(s,\delta)/\delta,w(s,\delta)/\delta)}.
\]
It has the required smooth one-sided extension with derivative $1$.
Symmetry gives the end at $\delta=L$ with derivative $-1$.
Away from the diagonal, smoothness is part of simplicity.
First variation and the zero diagonal value identify
\eqref{eq:primitive} with $D_g$, giving all of (C1).

The incoming flux measure is
\[
 d\mu=\sin\theta\dd s\dd\theta=\dd s\,|dp|.
\]
Changing from $p$ to the exit point with $p=-(D_g)_s$ and
\eqref{eq:twist} gives
\[
 \int_{\partial_+SM}\tau\dd\mu
     =\int_{\Cpair}D_g(D_g)_{st}\dd s\dd t.
\]
Santal\'o's formula \cite[Proposition~3.6.6]{PSU} says, for a continuous
base function $f$,
\begin{equation}\label{eq:santalo}
 \int_{\partial_+SM}\int_0^\tau f(c(t))\dd t\dd\mu
       =2\pi\int_M f\dd V_g.
\end{equation}
The factor $2\pi$ is the mass of the unit tangent circle in Liouville
measure $dV_g\,d\theta$ \cite[Section~3.5.2]{PSU}.
Taking $f=1$ proves \eqref{eq:santalo-area}.
\end{proof}

\subsection{Holomorphic traces determined by the boundary lift}
Let $X$ denote the geodesic vector field on $SM$, and $V$ the positive
fiber rotation. Put $X_\perp=[X,V]$ and
$\eta_\pm=(X\pm iX_\perp)/2$. A fiber mode $u_k$ satisfies
$-iVu_k=ku_k$; write
$\Omega_k=\{u\in C^\infty(SM;\mathbb C):-iVu=ku\}$.
Smooth fiber Fourier series and their projections converge
in $C^\infty$, and
\begin{equation}\label{eq:mode-operators}
 X=\eta_++\eta_-,\qquad
 (Xw)_k=\eta_+w_{k-1}+\eta_-w_{k+1},
\end{equation}
where $\eta_+$ raises and $\eta_-$ lowers the index
\cite[Lemmas~6.1.3 and 6.1.5]{PSU}. On base functions,
$\eta_-H=0$ is exactly holomorphicity in the oriented complex structure
\cite[Lemma~6.1.21]{PSU}. Vanishing of a mode and the zeroth average
are independent of an orthonormal frame rotation. In particular the
boundary tangent frame need not extend to the interior.

We use three transport facts with their smooth domains. On a simple
surface, a smooth function on $\partial SM$ invariant under the full
scattering involution extends uniquely to a smooth $w$ on $SM$ with
$Xw=0$ \cite[Theorem~5.1.1]{PSU}. On such a surface $I_0^*$ is
surjective onto $C^\infty(M)$; hence every smooth complex base function
$H$ is the zeroth Fourier coefficient of a smooth invariant function
\cite[Theorems~8.2.1 and 8.2.2]{PSU}. Here the source's
$\ell_0^*w=2\pi w_0$ is the unnormalized fiber integral, so the
source theorem is applied to $2\pi H$. Finally, on a compact nontrapping
surface with strictly convex boundary and surjective $I_0^*$, the ray
transform is injective on each single fiber mode: if $u$ is smooth,
$u|_{\partial SM}=0$, and $Xu$ has one mode, then $u=0$
\cite[Theorem~10.2.3]{PSU}. Indeed the integral of $Xu$ along every
geodesic is zero; injectivity gives $Xu=0$ and the zero boundary value
then gives $u=0$. Real statements extend to complex functions by applying
them to real and imaginary parts in the surjectivity step.

\begin{lemma}\label{lem:traces}
For actual simple-metric data, $\BA$ is exactly the set of boundary
restrictions of smooth holomorphic base functions on $(M,g)$.
\end{lemma}
\begin{proof}
Let $H$ be such a holomorphic function. Choose a smooth invariant $w$
with $w_0=H$ by the preceding surjectivity result, and replace $w$ by
$[w(x,v)+w(x,-v)]/2$. It remains invariant, since direction reversal
changes the sign of $X$, and its zeroth average remains $H$.
All its odd modes vanish. Put $W=\sum_{k\ge0}w_k$.
The projection is smooth through the boundary. In \eqref{eq:mode-operators},
every positive mode of $XW$ vanishes by $Xw=0$, its zeroth mode vanishes
by evenness, and its only possible negative mode is $\eta_-w_0=0$.
Thus $XW=0$. Its boundary trace $U$ is even, has no negative modes,
and is invariant under the full lift. Its average is the trace of $H$.

Conversely let $U$ satisfy \eqref{eq:boundary-fourier}. Its invariant
extension $w$ is smooth by the transport regularity statement.
Uniqueness follows by tracing an interior orbit to the incoming boundary.
The reversed $w(x,-v)$ has the same full boundary trace, by the evenness
of $U$, so uniqueness makes $w$ even. Set $w^- =\sum_{k<0}w_k$.
It is smooth with zero full boundary trace. The mode equations and
evenness give
\[
 Xw^-=\eta_+w_{-2}\in\Omega_{-1}.
\]
The single-mode injectivity statement applies with $m=-1$ and yields
$w^-=0$. The mode $-1$ equation for $Xw$ is now
$\eta_-w_0=0$. Hence $w_0$ is holomorphic and its boundary trace is
$U_0$. This proves the reverse inclusion.
\end{proof}

\subsection{Three-point normalization}\label{subsec:normalization}
\begin{lemma}\label{lem:coordinate}
For actual simple-metric data, (C2) holds. Its selected $\phi$ is the
boundary restriction of the unique holomorphic disk diffeomorphism with
the prescribed three-point normalization. In this coordinate the metric
is $\rho^2|dz|^2$ with $\rho$ smooth and positive on $\overline\D$ and
$\rho|_{\Sone}=b$.
\end{lemma}
\begin{proof}
Smooth uniformization \cite[Theorem~3.4.15]{PSU} supplies a holomorphic
map $\Phi_0:M\to\overline\D$ that is a smooth diffeomorphism of the
closed disks. Its boundary trace belongs to $\BA$ by
Lemma~\ref{lem:traces}.

We first show that every smooth holomorphic $H$ whose boundary trace
is an orientation-preserving circle diffeomorphism is itself a disk
diffeomorphism. In the coordinate $\Phi_0$, write $F=H\circ\Phi_0^{-1}$.
It is holomorphic on $\D$, smooth on its closure and nonconstant.
For a fixed $w$ with $|w|\ne1$, the functions $F-w$ are nonzero near
the boundary. The argument principle on circles $|z|=r<1$, close to $1$,
gives the zero count
\[
 n_r(w)=\frac1{2\pi i}\int_{|z|=r}\frac{F'(z)}{F(z)-w}\dd z.
\]
As $r\uparrow1$ this tends to the corresponding boundary integral.
Writing $F(e^{i\theta})=e^{ia(\theta)}$, where $a'>0$ and
$a(\theta+2\pi)=a(\theta)+2\pi$, and changing variable to $a$ shows
that the limit is $1$ for $|w|<1$ and $0$ for $|w|>1$.
For example these values follow by expanding the denominator in its
uniformly convergent geometric series on the boundary circle.
The integer $n_r(w)$ therefore equals its limit for all $r$ sufficiently
close to $1$. In the first case there are no zeros in a thin outer
annulus, so there is exactly one preimage in $\D$, with multiplicity one.
In the second case there is none: any preimage would lie inside one of
the circles with zero count.

An interior preimage of $w_0\in\Sone$ is also impossible. A small circle
around such a point has a positive argument-principle count for $F-w_0$.
That integer remains unchanged for $w$ near $w_0$, by continuity and
nonvanishing on the small contour. Taking $|w|>1$ contradicts the zero
count just proved. Thus $F$ maps $\D$ bijectively onto $\D$ and its
derivative never vanishes there. At the boundary its tangential derivative
is nonzero. The Cauchy--Riemann equations hold there by continuity, so
its real Jacobian is nonzero as well. A smooth extension and the inverse
function theorem give smooth local inverses through the boundary.
All uses of the argument principle were on interior contours; no analytic
continuation through $\Sone$ was assumed.

For completeness, the three-point normalization can be made explicitly.
For a positively ordered triple $\zeta_1,\zeta_2,\zeta_3$ on $\Sone$, set
\[
 C_{\zeta_3}(z)=i\frac{\zeta_3+z}{\zeta_3-z},\qquad
 r_j=C_{\zeta_3}(\zeta_j)\ (j=1,2),\qquad
 H_\zeta=\frac{C_{\zeta_3}-r_1}{r_2-r_1}.
\]
The imaginary part of $C_{\zeta_3}$ is
$(1-|z|^2)/|\zeta_3-z|^2$, and $r_1<r_2$ are real by the cyclic order.
Thus $H_\zeta$ maps the disk to the upper half-plane and the triple to
$(0,1,\infty)$. Applying the same formula to $(1,\omega,\omega^2)$,
the composition $H_{\rm roots}^{-1}H_\zeta$ is the desired disk
automorphism. Upper half-plane automorphisms are real fractional linear
maps \cite[Theorem~1.7]{Milnor}, and one fixing $0,1,\infty$ is the
identity. Any competing $\phi\in\BA$ extends by Lemma~\ref{lem:traces}
to a holomorphic $H$, which is a disk diffeomorphism by the preceding
argument. Comparing the two disk maps therefore proves uniqueness.

In the normalized coordinate the pushed metric is
$\rho^2|dz|^2$, with smooth positive $\rho$ through the boundary.
The boundary arclength identity is
\[
 \rho(\phi(s))|\phi_s(s)|=1,
\]
which gives exactly \eqref{eq:boundary-density}.
Together with Lemma~\ref{lem:geometry}, this proves (C2).
\end{proof}

\section{Variational realization}\label{sec:variation}
We now begin with proposed data, without assuming any realizing metric.
The first lemma makes the convex construction well defined on every
input passing its boundary prerequisites.

\subsection{Universal feasibility and projection}\label{subsec:projection}
\begin{lemma}\label{lem:projection}
Under (C1)--(C2), $\Fset$ is nonempty and convex. Its strong $H$ closure
$\Kset$ has a unique element $\rho_D$ of smallest norm. If this class
has a continuous representative on $\overline\D$, that representative
is unique.
\end{lemma}
\begin{proof}
Since $|D_t|=|Q|<1$ and the profile vanishes at both ends,
\[
 D(s,s+\delta)\le\delta,\qquad
 D(s,s+\delta)\le L-\delta.
\]
Thus $D$ is bounded by the shorter $h$-arclength arc. Put
$B=\max_{\Sone}b$. If $a,c\in\Sone$ have shorter angular separation
$\alpha\in(0,\pi]$, the corresponding $h$-shorter arc is no longer than
$B\alpha$. Since $\alpha\le(\pi/2)|a-c|$,
\[
 D_\phi(a,c)\le\frac\pi2B|a-c|.
\]
Every joining piecewise $C^1$ curve has Euclidean length at least
$|a-c|$, so the constant $\sigma=(\pi/2)B$ belongs to $\Fset$.

Strict positivity and the linear inequalities \eqref{eq:all-curves}
persist under convex combinations. Hence $\Fset$, and also its closure,
are convex. Let $m=\inf_{u\in\Kset}\|u\|_H$. Choose $u_n\in\Kset$
with $\|u_n\|_H^2\le m^2+1/n$. The midpoint is in $\Kset$, and
the parallelogram identity gives
\[
 \|u_n-u_k\|_H^2
 =2\|u_n\|_H^2+2\|u_k\|_H^2
       -4\|(u_n+u_k)/2\|_H^2\le 2/n+2/k.
\]
Completeness of $H$ and closedness of $\Kset$ produce a limit of norm
$m$. Applying the same identity to two minima proves uniqueness.
If continuous representatives differed at an interior point, their
difference would be bounded away from zero on a ball of positive area,
contradicting equality almost everywhere. Interior limits then identify
their boundary values.
\end{proof}

\subsection{The ODE tests and necessity}\label{subsec:necessity}
\begin{lemma}\label{lem:ode}
For a smooth positive $\rho$ on $\overline\D$, condition (C4) is
equivalent to simplicity of $g_\rho=\rho^2|dz|^2$.
\end{lemma}
\begin{proof}
The conformal geodesic equation has the first two components in
\eqref{eq:fixed-ode} \cite[Lemma~3.5.6]{PSU}. Its speed has
$g_\rho$ norm one. More explicitly the Christoffel symbols are
\[
 \Gamma^k_{ij}=\delta^k_i\lambda_j+\delta^k_j\lambda_i
                    -\delta_{ij}\lambda_k.
\]
In the curvature expression the quadratic first-derivative terms cancel,
giving $K_{g_\rho}=-\rho^{-2}(\lambda_{xx}+\lambda_{yy})$.
On the unit circle the unit tangent and inward normal are respectively
$\rho^{-1}(-\sin\psi,\cos\psi)$ and
$-\rho^{-1}(\cos\psi,\sin\psi)$. Substitution gives
\[
 g_\rho(\nabla_TT,\nu)=\rho^{-1}(1+\partial_r\lambda).
\]
Thus \eqref{eq:convexity} is exactly strict boundary convexity.
The third equation in \eqref{eq:fixed-ode} is the normal Jacobi equation
\cite[Lemmas~3.7.6 and 3.7.7]{PSU}.

Finite exit for all interior and inward boundary initial states gives
forward nontrapping; reversing the initial directions gives backward
nontrapping. A normal Jacobi field vanishing at its initial point is a
multiple of $j$. Its positivity through exit prevents conjugate pairs.
The tangential Jacobi component has the form $a+bt$ and cannot vanish
at two distinct times unless it is zero. Every subsegment begins at a
state covered by (C4), so this includes conjugate pairs starting in the
interior and pairs involving endpoints. The equivalence in
\cite[Theorem~3.8.2]{PSU} now gives simplicity. Conversely simplicity
gives finite exit and no zero of $j$ after its initial positive slope,
as well as strict convexity, so all the ODE tests hold.
\end{proof}

\begin{proposition}\label{prop:necessity}
Every smooth simple Riemannian metric inducing $h$ produces data
satisfying (C1)--(C5). Its normalized conformal density is exactly
$\rho_D$.
\end{proposition}
\begin{proof}
Lemmas~\ref{lem:geometry} and \ref{lem:coordinate} give (C1)--(C2)
and a normalized smooth positive density $\rho$ with boundary value $b$.
Every curve has $g_\rho$ length at least the boundary distance, so
$\rho\in\Fset$. For any $\sigma\in\Fset$ and any incoming
$g_\rho$-unit-speed geodesic, \eqref{eq:all-curves} gives
\[
 I_g(\sigma/\rho):=\int_0^\tau(\sigma/\rho)(c(t))\dd t\ge\tau.
\]
The endpoints are distinct. Integrate over the incoming boundary and
apply \eqref{eq:santalo} to the continuous base functions $\sigma/\rho$
and $1$. In Euclidean coordinates this yields
\[
 2\pi\int_\D\sigma\rho\dd x\dd y
        \ge 2\pi\int_\D\rho^2\dd x\dd y.
\]
The pairing with $\rho$ is continuous in $H$, so
$\langle u,\rho\rangle_H\ge\|\rho\|_H^2$ for every $u\in\Kset$.
Consequently
\[
 \|u\|_H^2
 =\|\rho\|_H^2+\|u-\rho\|_H^2+2\langle u-\rho,\rho\rangle_H
 \ge\|\rho\|_H^2+\|u-\rho\|_H^2.
\]
Thus the true density is the unique least-norm element $\rho_D$.
It satisfies (C3), and Lemma~\ref{lem:ode} gives (C4).
Finally \eqref{eq:santalo-area} gives
$\|\rho_D\|_H^2=\Area_g(M)=\Varea(D)$, which is (C5).
\end{proof}

\subsection{Passing length inequalities through the closure}
\begin{lemma}\label{lem:closure}
Assume (C1)--(C4), and use $\phi$ to express the boundary distance of
the selected simple metric $g_\rho$ in the original arclength
coordinates. Then
\begin{equation}\label{eq:distance-order}
 D_{g_\rho}(s,t)\ge D(s,t)\quad(s\ne t).
\end{equation}
\end{lemma}
\begin{proof}
The selected density has the correct boundary metric since
$\rho(\phi(s))|\phi_s(s)|=1$. Choose $\sigma_n\in\Fset$ with
$\sigma_n\to\rho$ strongly in $H$, and put
$f_n=(\sigma_n-\rho)/\rho$. All these functions are smooth.
Santal\'o's formula and Cauchy--Schwarz give
\begin{align}\label{eq:family-L1}
 \int_{\partial_+S\overline\D}|I_gf_n|\dd\mu
 &\le\int_{\partial_+S\overline\D}I_g|f_n|\dd\mu\\
 &=2\pi\int_\D|\sigma_n-\rho|\rho\dd x\dd y
 \le2\pi\|\sigma_n-\rho\|_H\|\rho\|_H\longrightarrow0.\notag
\end{align}
For each of its geodesics, the curve inequality gives
$\tau+I_gf_n\ge D(s,t)$. Therefore
\[
 (D(s,t)-\tau)_+\le |I_gf_n|.
\]
The integral of this nonnegative left side is zero by
\eqref{eq:family-L1}. It is continuous on the open incoming bundle,
where flux has strictly positive density. Positivity at one state would
give a neighborhood of positive integral; hence it vanishes everywhere.
The endpoint diffeomorphism of Lemma~\ref{lem:geometry} then gives
\eqref{eq:distance-order} for every distinct boundary pair.
Only smooth functions were integrated along individual curves. No trace
on a curve or ray transform of a general $L^2$ class was used.
\end{proof}

\subsection{Strict area monotonicity and sufficiency}
\begin{lemma}\label{lem:area-order}
Let $D_0,D_1$ have smooth zero-ended profiles as in
\eqref{eq:profile}, and suppose $(D_i)_{st}>0$ on $\Cpair$.
If $D_1\ge D_0$ and $\Varea(D_1)=\Varea(D_0)$, then $D_1=D_0$.
\end{lemma}
\begin{proof}
On the compact profile rectangle put
$\mathcal L=\partial_s\partial_\delta-\partial_\delta^2$.
For smooth periodic zero-ended profiles $u,v$, two integrations by parts
give
\[
 \int u v_{s\delta}=-\int u_\delta v_s=\int u_{s\delta}v,
 \qquad
 -\int u v_{\delta\delta}=\int u_\delta v_\delta
                           =-\int u_{\delta\delta}v.
\]
Terms at $s=0,L$ cancel by periodicity, while the first and second
$\delta$ boundary terms vanish because $u=0$ and $v=0$, respectively.
Thus $\int u\mathcal Lv=\int v\mathcal Lu$. Expanding the quadratic
functional \eqref{eq:area-action} gives
\begin{equation}\label{eq:area-difference}
 \Varea(D_1)-\Varea(D_0)
 =\frac1{2\pi}\int_{\Cpair}(D_1-D_0)
             \big((D_1)_{st}+(D_0)_{st}\big)\dd s\dd t.
\end{equation}
If $D_1-D_0$ is nonzero and nonnegative, continuity makes it positive
on an open set. The mixed derivatives are positive there, so the right
side is strictly positive, a contradiction. All integrations by parts
occur on the smooth compact rectangle; no diagonal-strip limit is needed.
\end{proof}

\begin{proof}[Proof of Theorem~\ref{thm:main}]
Necessity is Proposition~\ref{prop:necessity}. For sufficiency assume
all five conditions. Lemma~\ref{lem:projection} defines the selected
density, (C3) supplies its smooth positive representative, and
Lemma~\ref{lem:ode} makes $g_\rho$ simple. Lemma~\ref{lem:closure}
gives $D_{g_\rho}\ge D$. The two actions have smooth zero-ended
profiles and positive mixed derivative, by (C1) and
Lemma~\ref{lem:geometry}. Their action areas agree, since
\[
 \Varea(D_{g_\rho})=\Area_{g_\rho}(\overline\D)
                 =\|\rho_D\|_H^2=\Varea(D).
\]
Lemma~\ref{lem:area-order} forces $D_{g_\rho}=D$ pointwise.

For $(s,p)\in Z$ let $A(s,p)=(t,q)$. The unique geodesic from $s$
to $t$ for the selected metric has initial tangential covector
$-(D_{g_\rho})_s=-D_s=P(s,t)=p$ and exit covector
$(D_{g_\rho})_t=D_t=Q(s,t)=q$. Thus its scattering is exactly $A$
on the whole open cylinder.

To realize this metric on the fixed smooth disk $M$, choose any
orientation-preserving smooth disk chart $M\to\overline\D$.
In this chart the desired boundary map is
$e^{i\theta}\mapsto e^{ia(\theta)}$ with $a'>0$ and
$a(\theta+2\pi)=a(\theta)+2\pi$. Choose a smooth radial cutoff
$0\le\chi\le1$, zero near the origin and one near the boundary.
The map
\[
 (r,\theta)\longmapsto
       (r,\theta+\chi(r)(a(\theta)-\theta))
\]
is well defined modulo $2\pi$ and has positive angular derivative
$1-\chi+\chi a'$. It is the identity near the center and a smooth
bijection on each circle, with smooth inverse. Composing with the chart
gives a smooth disk diffeomorphism $\Psi$ with boundary value $\phi$.
The pullback $\Psi^*(\rho^2|dz|^2)$ is a smooth simple Riemannian
metric on $M$. It induces $h$ by \eqref{eq:rho-trace} and
\eqref{eq:boundary-density}, and its scattering is the prescribed $A$.

Every initial choice was arbitrary in both directions of the proof.
Passing the tests for one choice gives a realizing metric; necessity
then makes its data pass for every other choice. This proves the
claimed independence of the classification. For proposed data without
a realization, no algebra or fractional-linear closure property of
$\BA$ has been assumed.
\end{proof}

\paragraph{AI-use statement.}
An OpenAI Codex agent assisted in developing the argument, checking cited
source statements, and preparing this manuscript. The submitted result is
an analytic proof claim. AI-assisted checks and successful typesetting do
not establish independent mathematical verification. Independent review
of the complete argument is requested.

\begin{thebibliography}{9}
\bibitem{BI}
D. Burago and S. Ivanov,
\emph{Boundary distance, lens maps and entropy of geodesic flows of Finsler metrics},
Geometry \& Topology \textbf{20} (2016), 469--490.
\href{https://doi.org/10.2140/gt.2016.20.469}{doi:10.2140/gt.2016.20.469}.

\bibitem{BP}
J. Bohr and G. P. Paternain,
\emph{The Transport Oka--Grauert Principle for Simple Surfaces},
Journal de l'\'Ecole polytechnique --- Math\'ematiques \textbf{10} (2023), 727--769.
\href{https://doi.org/10.5802/jep.231}{doi:10.5802/jep.231}.

\bibitem{Milnor}
J. Milnor,
\emph{Dynamics in One Complex Variable: Introductory Lectures},
Stony Brook IMS Preprint \#1990/5, revised September 1991.
\url{https://abel.math.harvard.edu/archive/118r_spring_05/docs/milnor.pdf}.

\bibitem{PSU}
G. P. Paternain, M. Salo and G. Uhlmann,
\emph{Geometric Inverse Problems: With Emphasis on Two Dimensions},
Cambridge Studies in Advanced Mathematics, vol.~204,
Cambridge University Press, 2023.
\href{https://doi.org/10.1017/9781009039901}{doi:10.1017/9781009039901}.

\bibitem{PU}
L. Pestov and G. Uhlmann,
\emph{Two dimensional compact simple Riemannian manifolds are boundary distance rigid},
Annals of Mathematics (2) \textbf{161} (2005), 1093--1110.
\href{https://doi.org/10.4007/annals.2005.161.1093}{doi:10.4007/annals.2005.161.1093}.
\end{thebibliography}
\end{document}
